\documentclass[12pt,addpoints,answers]{repaso2}
\grado{2}
\nivel{Secundaria}
\cicloescolar{2026-2027}
\materia{Matemáticas}
\unidad{1}
\title{Repaso extra}
\aprendizajes{
      \item Resuelve problemas que impliquen la suma, resta, la multiplicación y la división de números enteros, aplicando las reglas correspondientes.
      \item Identifica y ubica números negativos en una recta numérica, comparando su magnitud.
      \item Aplica las propiedades de las potencias a números negativos en la resolución de problemas.
      \item Resuelve problemas que involucren las leyes de los exponentes, y expresa números en notación científica.
      \item Resuelve problemas de contexto científico y tecnológico utilizando la notación científica.
      \item Identifica y ubica puntos en el plano cartesiano, y comprende la estructura de los cuadrantes.
      \item Calcula la pendiente de una recta y comprende su significado en diferentes contextos.
      \item Resuelve problemas que involucren el uso de porcentajes, mediante la ``regla de tres''.
}
\author{Melchor Pinto, J.C.}
\begin{document}
\INFO
\newpage
\begin{multicols}{2}
	\tableofcontents
\end{multicols}
\newpage
\begin{questions}
		\section{Unidad 1}
	\subsection{Cálculos numéricos}
	%\subsubsection{Suma de números}
	 \questionboxed[1]{}
	% \subsubsection{Resta de números}
	% % \questionboxed[1]{}
	% \subsubsection{Multiplicación de números}
	% % \questionboxed[1]{}
	% \subsubsection{División de números}
	% % \questionboxed[1]{}
	% \subsubsection{Resolución de Problemas}
	% % \questionboxed[1]{}
	% \subsection   {Fracciones}
	% \subsubsection{Clasificación de fracciones}
	% % \questionboxed[1]{}
	% \subsubsection{Representación de fracciones}
	% % \questionboxed[1]{}
	% \subsubsection{Nombre de fracciones}
	% % \questionboxed[1]{}
	% \subsubsection{Fracciones en la recta numérica}
	% % \questionboxed[1]{}
	% \subsubsection{Conversión de fracciones}
	% % \questionboxed[1]{}

	% \subsection   {Fracciones MCD y MCM}
	% \subsubsection{Comparación de fracciones}
	% % \questionboxed[1]{}
	% \subsubsection{Fracciones equivalentes}
	% % \questionboxed[1]{}
	% \subsubsection{M.C.D y M.C.M}
	% % \questionboxed[1]{}
	% \subsubsection{Simplificación de fracciones}
	% % \questionboxed[1]{}
	% \subsubsection{Resolución de problemas}
	% % \questionboxed[1]{}

	% \subsection   {Números decimales}
	% \subsubsection{Ubicación en la recta numérica}
	% % \questionboxed[1]{}
	% \subsubsection{Porcentajes a decimal}
	% % \questionboxed[1]{}
	% \subsubsection{Operaciones con múltiplos de 10}
	% % \questionboxed[1]{}
	% \subsubsection{Conversión de fracciones a decimales}
	% % \questionboxed[1]{}
	% \subsubsection{Conversión de decimales a fracciones}
	% % \questionboxed[1]{}

	% \subsection   {Conceptos básicos de números negativos}
	% \subsubsection{Ubicación en la recta numérica}
	% % \questionboxed[1]{}
	% \subsubsection{Comparación de negativos}
	% % \questionboxed[1]{}
	% \subsubsection{Oraciones con negativos}
	% % \questionboxed[1]{}
	% \subsubsection{Determina el signo en suma y resta con negativos}
	% % \questionboxed[1]{}
	% \subsubsection{Determina el signo multiplicación y división con negativos}
	% % \questionboxed[1]{}

	% \subsection   {Operaciones con números negativos}
	% \subsubsection{Comparación de negativos}
	% % \questionboxed[1]{}
	% \subsubsection{Suma y resta con negativos}
	% % \questionboxed[1]{}
	% \subsubsection{Multiplicación y división con negativos}
	% % \questionboxed[1]{}
	% \subsubsection{Potencias con numeros negativos}
	% % \questionboxed[1]{}
	% \subsubsection{Jerarquía de operacione}
	% % \questionboxed[1]{}

	% \section{Unidad 2}
	% \subsection   {Círculo}
	% \subsubsection{Resolución de problemas}
	% % \questionboxed[1]{}
	% \subsubsection{Radio, Diámetro, Perímetro y Área de un círculo}
	% % \questionboxed[1]{}
	% \subsection   {Polígonos y circunferencias}
	% \subsubsection{Ángulos interiores}
	% % \questionboxed[1]{}
	% \subsubsection{Ángulos centrales y exteriores}
	% % \questionboxed[1]{}
	% \subsubsection{Ángulos centrales e inscritos}
	% % \questionboxed[1]{}
	% \subsubsection{Arco de una circunferencia}
	% % \questionboxed[1]{}
	% \subsubsection{Área de un sector circular}
	% % \questionboxed[1]{}
	% \subsection   {Figuras y cuerpos geométricos}
	% \subsubsection{Perímetro y Área}
	% % \questionboxed[1]{}
	% \subsubsection{Resolución de problemas}
	% % \questionboxed[1]{}
	% \subsubsection{Área lateral, Área total y Volumen}
	% % \questionboxed[1]{}
	% \subsection   {Monomios y polinomios}
	% \subsubsection{Lenguaje algebraico}
	% % \questionboxed[1]{}
	% \subsubsection{Suma de monomios y polinomios}
	% % \questionboxed[1]{}
	% \subsubsection{Resta de monomios y polinomios}
	% % \questionboxed[1]{}
	% \subsubsection{Operaciones combinadas}
	% % \questionboxed[1]{}
	% \subsubsection{Perímetro de figuras geométricas}
	% % \questionboxed[1]{}
	% \subsection   {Operaciones con monomios y polinomios}
	% \subsubsection{Suma, resta y multiplicación de exponentes}
	% % \questionboxed[1]{}
	% \subsubsection{Multiplicación y división de monomios y polinomios}
	% % \questionboxed[1]{}
	% \subsubsection{Áreas de figuras geométricas}
	% % \questionboxed[1]{}
	% \subsection   {Sistema de unidades}
	% \subsubsection{Unidades de longitud}
	% % \questionboxed[1]{}
	% \subsubsection{Unidades de masa}
	% % \questionboxed[1]{}
	% \subsubsection{Unidades de capacidad}
	% % \questionboxed[1]{}
	% \subsubsection{Unidades de área y volumen}
	% % \questionboxed[1]{}

	% \section{Unidad 3}
	% \subsection   {Probabilidad y estadística}
	% \subsubsection{Mediana y moda}
	% % \questionboxed[1]{}
	% \subsubsection{Promedio}
	% % \questionboxed[1]{}
	% \subsubsection{Interpretación de gráficas}
	% % \questionboxed[1]{}
	% \subsubsection{Eventos mutuamente excluyentes}
	% % \questionboxed[1]{}
	% \subsubsection{Eventos dependientes e independientes}
	% % \questionboxed[1]{}
	% \subsection   {Razones y proporciones}
	% \subsubsection{Relaciones proporcionales}
	% % \questionboxed[1]{}
	% \subsubsection{Constante de proporcionalidad}
	% % \questionboxed[1]{}
	% \subsubsection{Regla de correspondencia (ecuación)}
	% % \questionboxed[1]{}
	% \subsubsection{Proporción directa e inversa}
	% % \questionboxed[1]{}
	% \subsubsection{Proporciones compuestas}
	% % \questionboxed[1]{}
	% \subsection   {Sucesiones aritméticas}
	% \subsubsection{Completando la sucesión}
	% % \questionboxed[1]{}
	% \subsubsection{Diferencia de una sucesión}
	% % \questionboxed[1]{}
	% \subsubsection{Término enésimo}
	% % \questionboxed[1]{}
	% \subsubsection{Término general}
	% % \questionboxed[1]{}
	% \subsubsection{Término enésimo 2}
	% % \questionboxed[1]{}
	% \subsection   {Ecuaciones lineales}
	% \subsubsection{Lenguaje algebraico}
	% % \questionboxed[1]{}
	% \subsubsection{Sustitución de valores}
	% % \questionboxed[1]{}
	% \subsubsection{Ecuaciones de primer grado}
	% % \questionboxed[1]{}
	% \subsubsection{Resolución de problemas}
	% % \questionboxed[1]{}
	% \subsection   {Sistemas de ecuaciones}
	% \subsubsection{Método de eliminación}
	% % \questionboxed[1]{}
	% \subsubsection{Método de sustitución}
	% % \questionboxed[1]{}
	% \subsubsection{Método de igualación}
	% % \questionboxed[1]{}
	% \subsubsection{Sistema de ecuaciones}
	% % \questionboxed[1]{}

\end{questions}
\end{document}